%% ============================================================================
%%  main.tex -- Esqueleto de armado del informe
%%
%%  Este archivo NO contiene texto de la propuesta: solo metadatos y el orden
%%  de los subdocumentos. Todo el contenido vive en carpetas, una por
%%  subdocumento del Formulario T-7.
%%
%%  Compilar:  latexmk main.tex
%%  Limpiar:   latexmk -c
%%
%%  COMO AGREGAR UN SUBDOCUMENTO
%%  ----------------------------
%%  1. Deja la carpeta al lado de este archivo (o en cualquier ruta relativa).
%%  2. Dentro debe haber un contenido.tex que empiece por \chapter{...}.
%%  3. Agrega una linea \subdocumento{nombre_de_la_carpeta} en el orden que
%%     corresponda. Eso es todo: no hay que tocar nada mas.
%%
%%  La carpeta es autocontenida. Sus figuras van en figuras/ y se llaman por
%%  ruta relativa (figuras/edt.png); sus fragmentos, con \parte{tablas/x}.
%%  Nada dentro del capitulo sabe donde esta montado, asi que reubicarlo es
%%  mover la carpeta y cambiar la ruta de esta llamada.
%% ============================================================================

% Opciones: borrador | final | firma | sinfirma | carta | oficio | apa | sinnotas
\documentclass[carta,firma,borrador]{lafrox}

% ---------------------------------------------------------------------------
%  Metadatos de la oferta. Se editan SOLO aqui: caratula, encabezados, pies y
%  metadatos del PDF se alimentan de estas claves. Las que no se declaran
%  toman el valor por defecto de lafrox.cls (seccion 6).
% ---------------------------------------------------------------------------
\datoslafrox{
  documento  = Propuesta Técnica,
  subtitulo  = Informe 2,
  alcance    = Subdocumentos 1--9 y 13,
  formulario = Formulario T-22,
  fecha      = 05 de octubre de 2026,
  version    = 1.0,
}

\begin{document}

\portadalafrox
\indicelafrox

% ---------------------------------------------------------------------------
%  SUBDOCUMENTOS DEL FORMULARIO T-7
%
%  Descomenta cada linea a medida que su carpeta exista. El orden de estas
%  llamadas ES el orden del informe; la numeracion de capitulos, figuras y
%  tablas se recalcula sola.
%
%  El Informe 2 (T-22, 05-10-2026) cubre los subdocumentos 1 a 9 y el 13.
% ---------------------------------------------------------------------------

\subdocumento{subdoc}   % <- borrar al montar el informe real
\subdocumento{subdocumento-ejemplo}
% \subdocumento{01_presentacion_empresa}
% \subdocumento{02_problema_necesidad}
% \subdocumento{03_esquema_solucion_alcance}
% \subdocumento{04_arquitectura}
% \subdocumento{05_modelo_datos}
% \subdocumento{06_metodologias}
% \subdocumento{07_plan_trabajo_edt}
% \subdocumento{08_plan_riesgos}
% \subdocumento{09_plan_calidad}
% \subdocumento{13_innovaciones}

%  Subdocumentos 10, 11, 12 y 14: entran en el Informe 3.
% \subdocumento{10_operacion_niveles_servicio}
% \subdocumento{11_planes_operacion}
% \subdocumento{12_equipo_subcontrataciones}
% \subdocumento{14_ventajas_beneficios_consolidacion}

\end{document}
